\section{A whole-source prime bound for the elliptic Heisenberg escape}
\label{sec:elliptic-heisenberg-prime}

\paragraph{Status and provenance.}
D-TD55, question Q2fk, is a separate consequence of the exact
projective Heisenberg trace and genus results D-TD05--10.  The parent
asked whether a uniform support floor for the entire original image
could control the remaining elliptic escape.  The argument below was
developed by \texttt{/root/imprimitive}; the parent independently checked
the proposed budget inequalities.  The complete author proof, source,
checker and first-output readings pass.  Complete root and independent
whole-proof reviews also pass.  The result does not bound the
Heisenberg exponent at the remaining small primes.

\paragraph{Actual source and the full-image hypothesis.}
Let \(C\to\mathbf P^1\) be a connected finite \(J\)-Galois cover, with
faithful irreducible complex coefficient representation on \(V\),
\(e=\dim V\).  Let \(T_1,\ldots,T_r\) be its generating product-one
inertia tuple, with every \(T_i\ne I\), and suppose
\[
 r\ge12,\qquad
 \sum_{i=1}^r\operatorname{rank}(T_i-I)=2e+6,\qquad
 \frac{\operatorname{rank}(T_i-I)}e\ge\frac1M
 \quad(1\le i\le r),
 \qquad M\ge1.
 \tag{EHP.1}
\]
The last condition is on the actual original tuple.  One sufficient
condition is the projective support floor
\(\nu_V(T)/e\ge1/M\) for every nonscalar element of \(J\), since a
nonidentity scalar has moving rank \(e\).  No such floor is inferred
from only one selected tensor factor.

Suppose \(V=W\otimes U\) is preserved factorwise by \(J\).  Reciprocal
scalar ambiguities are allowed.  Its actual finite projective image
\(P\subset\operatorname{PGL}(W)\) contains a normal subgroup
\[
 Q\simeq\mathbf F_p^{\,2m},\qquad m\ge1,
 \qquad
 [L_v,L_w]=\zeta_p^{\langle v,w\rangle}I,
 \tag{EHP.2}
\]
with nondegenerate alternating pairing, and the projective action of
\(Q\) on \(W\) is irreducible.  Write
\[
 Z=C/\ker(J\to P),\qquad Z/P=\mathbf P^1,\qquad
 y_i=\pi_W(T_i),\qquad
 k=\#\{i:y_i\ne1\}.
 \tag{EHP.3}
\]
Here \(\pi_W:J\to P\) is the actual projective factor homomorphism;
no linear restriction of \(T_i\) to \(W\) is being asserted.

\begin{theorem}[D-TD55]
In the setup above, if \(p\ge5\), \(m=1\), and \(g(Z)=1\), then
\[
 p\le\frac{8M}{r-4}\le M.
 \tag{EHP.4}
\]
If the projective factor \(P\) is primitive, then the stronger bound is
\[
 p\le\frac{8M}{r-3}\le\frac{8M}{9}.
 \tag{EHP.5}
\]
Consequently every source satisfying (EHP.1)--(EHP.3), with arbitrary
\(m\), satisfies
\[
 p\le\max\!\left\{7743,\frac{8M}{r-4}\right\}
       \le\max\{7743,M\}.
 \tag{EHP.6}
\]
If \(P\) is primitive, replace \(r-4\) by \(r-3\) in (EHP.6).
In particular that replacement is valid when the original coefficient
group \(J\) is primitive.
\end{theorem}

\begin{proof}
\emph{Step 1: the actual elliptic quotient and its stabilizers.}
Assume first \(p\ge5\), \(m=1\), and \(g(Z)=1\).  Choose an origin on
\(Z\).  The origin-fixing automorphism group is cyclic of order two,
four, or six, and all automorphisms are translations followed by such
an automorphism.  These are exactly
\cite[Chapter IV, Proposition 7.13 and Remark 7.18, pp.~140--141]{TD:MilneEC2021}.
Since \(p\ge5\), the elementary abelian \(p\)-group \(Q\) has trivial
image in that cyclic quotient.  Thus \(Q\) consists of translations.
The translation subgroup of \(P\) centralizes \(Q\) projectively.
The exact centralizer conclusion D-TD05 therefore makes that subgroup
equal to \(Q=Z[p]\).

It follows that \(P/Q\) is cyclic of order \(d\in\{1,2,3,4,6\}\).
The quotient of an elliptic curve by a finite translation group is
elliptic, so \(Z/P=\mathbf P^1\) excludes \(d=1\).  An automorphism
with nontrivial linear part has a fixed point: on the complex torus
it is \(z\mapsto\lambda z+b\), and \(1-\lambda\) is a nonzero,
surjective isogeny.  Taking a fixed point of a generator as origin
gives
\[
                 P=Z[p]\rtimes C_d,\qquad d\in\{2,3,4,6\}.
 \tag{EHP.7}
\]
This is the already proved D-TD07 argument, whose use of the complex
torus is covered by
\cite[Chapter III, Proposition 3.3 and Theorem 3.10,
pp.~94, 96]{TD:MilneEC2021}; it only needs \(p\ge5\) in the present
elliptic case.

Every nonidentity inertia image \(y_i\) fixes a point of \(Z\).
A nonidentity translation has no fixed point.  Moreover a point
stabilizer injects into the origin-fixing cyclic group: its kernel
would be a translation fixing that point.  Consequently, for its
exact projective order \(o_i\), the linear part of \(y_i\) has the
same order \(o_i\), with \(o_i\mid d\) and \(o_i\ge2\).
For \(0<j<o_i\), the linear part of \(y_i^j\) is still nonidentity.

Its action on \(Q=Z[p]\) is the reduction of the action on the
rank-two period lattice.  For a nonidentity linear part of order
\(2,3,4,6\), respectively, its characteristic polynomial is
\[
 (X+1)^2,\qquad X^2+X+1,\qquad X^2+1,\qquad X^2-X+1.
 \tag{EHP.8}
\]
The values at \(X=1\) are \(4,3,2,1\), respectively.  They are all
nonzero modulo \(p\ge5\).  Hence every proper nonidentity power of
each \(y_i\) acts on the label space \(Q\) with zero fixed space.
The nondegenerate Heisenberg degree and exact trace formula D-TD05
now give
\[
 \dim W=p,\qquad
 \frac{|\operatorname{tr}A_i^j|}{p}
 =
 \begin{cases}
 1,&o_i\mid j,\\
 1/p,&o_i\nmid j,
 \end{cases}
 \tag{EHP.9}
\]
for a unitary factor representative \(A_i\) and every integer \(j\).
In the second case the fixed-space phase sum consists of the single
label zero, so its value is exactly one.  In the first case the
factor is scalar.  No splitting of a general projective normalizer
is being assumed.

\emph{Step 2: every original period, including projected identity powers.}
Let \(L_i\) be the exact order of the original coefficient operator
\(T_i\).  The projective homomorphism implies \(o_i\mid L_i\).
Write \(T_i=A_i\otimes B_i\), with both factors unitary after
reciprocal scalar normalization.  Indeed \(T_i^{L_i}=I\) implies
\(A_i^{L_i}\) and \(B_i^{L_i}\) are reciprocal scalars; scaling one
\(L_i\)-th root makes both have finite order, and an invariant
Hermitian form then unitarizes each.  The absolute trace ratios
in (EHP.9) are unchanged by those scalars.

The \(1\)-eigenspace multiplicity of \(T_i\) is its ordinary full-period
Fourier average.  Using \(|\operatorname{tr}B_i^j|\le\dim U\) and
retaining all \(L_i/o_i\) powers whose projected action is the
identity gives
\[
 \begin{split}
 \frac{\dim\ker(T_i-I)}e
 &\le \frac1{L_i}\sum_{j=0}^{L_i-1}
                    \frac{|\operatorname{tr}T_i^j|}{e}\\
 &\le \frac1{o_i}+\frac{1-1/o_i}{p},\\
 \frac{\operatorname{rank}(T_i-I)}e
 &\ge (1-1/p)(1-1/o_i).
 \end{split}
 \tag{EHP.10}
\]
The same calculation works for every eigenvalue target.  Thus
arbitrary original scalars, extra scalar order, and the entire
tensor partner are all retained.  For \(y_i=1\) this estimate would
only give zero; those original entries instead retain their
actual lower bound \(1/M\) from (EHP.1).

\emph{Step 3: projected branching and the original missing marks.}
Riemann--Hurwitz for the connected actual \(P\)-cover \(Z\to\mathbf P^1\)
gives
\[
                  \sum_{y_i\ne1}(1-1/o_i)=2.
 \tag{EHP.11}
\]
Here each index is an original branch value with nonidentity image;
no projected branch values are introduced or identified.  Since
each summand is at least \(1/2\), \(k\le4\).
The nonidentity projected images still generate \(P\), because
the original tuple generates \(J\).

If \(P\) is primitive, \(k=4\) is impossible.  Equality in the
four-mark bound would give \(o_i=2\) for all four marks.  All their
nontrivial linear parts are the unique involution of the cyclic
group \(C_d\).  Since those images generate \(P/Q\), this forces
\(d=2\).  Its action on \(Q\) is \(-I\), preserving every line of
\(\mathbf F_p^2\).  Lifts of a line have \(p\) distinct one-dimensional
common eigenspaces on \(W\), and the full normalizer of that line
permutes these eigenspaces.  They give a nontrivial imprimitivity
system for \(P\), as proved explicitly in D-TD07.  Thus
\[
                k\le4,\qquad
                P\text{ primitive}\ \Longrightarrow\ k\le3.
 \tag{EHP.12}
\]

Summing (EHP.10) over the projected genuine marks and the actual
whole-image floor over the other \(r-k\) original marks yields
\[
 2+\frac6e
 \ge 2-\frac2p+\frac{r-k}{M},
 \qquad
 \frac{r-k}{M}\le\frac2p+\frac6e\le\frac8p,
 \tag{EHP.13}
\]
because \(e=p\dim U\ge p\).  Combining (EHP.12) and (EHP.13)
proves (EHP.4)--(EHP.5).

\emph{Step 4: all Heisenberg primes.}
If \(p<7744\), the integer bound \(p\le7743\) suffices.
If \(p\ge7744\), D-TD06 forces \(m=1\) and \(g(Z)=1\);
the preceding argument applies.  This proves (EHP.6).
Finally, an imprimitivity decomposition of \(W\) for \(P\) tensors
with the whole \(U\) to one of \(V\) for \(J\).
Therefore a primitive original coefficient image makes the actual
factor image primitive, exactly as in D-TD11.
\end{proof}

\paragraph{Conditional uniform application, with no cutoff presumed.}
Suppose a faithful primitive source as above has one complete D-TD12
recursion for which every alternating simple quotient is \(A_n\) with
\(n<N\), where \(N\ge82\) is an integer.  D-TD13 supplies the whole
original-image floor with
\[
             M_0=\max\{80,N-2\}=N-2.
 \tag{EHP.14}
\]
Every original nonidentity scalar has full moving rank, so (EHP.1)
holds with \(M=M_0\).  Apply the theorem to each actual stable
Heisenberg coarse factor supplied by that recursion, keeping its
entire complementary factor.  Each such factor is primitive by
D-TD11.  Its prime therefore satisfies
\[
                   p\le
           \max\!\left\{7743,\frac{8(N-2)}9\right\}.
 \tag{EHP.15}
\]
This implication does not assume that an alternating cutoff has
already been proved in every source class.  In particular substituting
\(N=N_*\) is justified only for the exact minimal-source comparison
class in which the separate alternating-cutoff synthesis is accepted.
The theorem itself is the explicit-\(M\) statement (EHP.4)--(EHP.6).

\paragraph{Remaining boundary.}
The proof uses at least eight original genuine marks which vanish in
the elliptic projective quotient, or at least nine in the primitive
case.  It does not infer twelve projected marks.  The elliptic
three-mark examples of D-TD07 therefore do not contradict this
whole-source restriction and are still not claimed to be complete
sixfold suppliers.  At primes below \(7744\), the exponent \(m\)
can still be arbitrarily large in the currently retained necessary
conditions.  No bound on that exponent, the coefficient dimension,
the full image order, original field descent, period-map image,
polarization realization, or a detecting cycle is asserted here.
